\chapter{Charakterystyka przedsięwzięcia}

\section{Cel projektu}

\textbf{E-Learning} to webowa platforma kursów online zbudowana w~jednym spójnym systemie
Django. Łączy cztery obszary, które w~typowych rozwiązaniach są rozproszone między osobnymi
narzędziami: \emph{sklep} z~kursami (koszyk, zamówienia, płatność online), \emph{środowisko
nauki} (lekcje, testy z~ocenianiem automatycznym i~ręcznym, dyplomy), \emph{moderowaną
społeczność} (komentarze w~wątkach, głosy, rangi, zgłaszanie kursów) oraz \emph{asystenta
AI} opartego o~polski model językowy Bielik, uruchamiany lokalnie przez Ollamę.

Projekt powstał w~ramach przedmiotu \emph{Inteligentne multimedialne systemy uczące}. Poza funkcjami biznesowymi
istotnym celem była \textbf{dostępność i~użyteczność}: cały serwis ma być obsługiwalny wyłącznie klawiaturą,
zgodny z~WCAG 2.1 na poziomie AA i~zbudowany bez gotowych bibliotek komponentów
(Bootstrap, Tailwind, szablony). Wygląd platformy opiera się na autorskim arkuszu stylów
sterowanym pięcioma kolorami bazowymi, które administrator może zmieniać z~panelu.

Niniejszy dokument opisuje system w~stanie po zakończeniu prac implementacyjnych. Zawiera
analizę wymagań, model aktorów i~przypadków użycia, architekturę, model danych, opis
interfejsów programistycznych, kluczowych algorytmów, zagadnień dostępności, przegląd
interfejsu użytkownika, sposób weryfikacji, przebieg realizacji oraz instrukcję
uruchomienia. Odbiorcami dokumentu są osoba oceniająca projekt oraz każdy, kto chciałby
system uruchomić lub rozwijać.

\section{Słownik pojęć}

\begin{table}[H]
\centering\small
\caption{Słownik pojęć używanych w~dokumentacji}
\begin{tabularx}{\linewidth}{@{}P{3.6cm}L@{}}
\toprule
\hd{Pojęcie} & \hd{Znaczenie} \\
\midrule
Kurs & Produkt w~katalogu: opis, cena (lub bezpłatny), poziom, kategoria, program złożony z~sekcji i~lekcji oraz testów. \\
Program kursu & Uporządkowana lista sekcji (\code{CourseSection}), w~każdej lekcje (\code{Lesson}); testy mogą stać po wybranej lekcji albo na końcu kursu. \\
Test & Zestaw pytań (\code{Quiz}): jednokrotnego wyboru (ABC), wielokrotnego wyboru i~otwartych. Test obowiązkowy blokuje dalszą naukę do czasu zaliczenia. \\
Podejście & Jedno wysłanie testu przez ucznia (\code{Attempt}) wraz z~odpowiedziami i~wynikiem. \\
Zapis & Powiązanie ucznia z~kursem (\code{Enrollment}), powstaje po zapisie na kurs bezpłatny, opłaceniu zamówienia albo wykorzystaniu kodu dostępu. \\
Zgłoszenie kursu & Kurs przesłany do akceptacji przez osobę bez uprawnień moderatora danej kategorii; ma trwały wątek dyskusji z~nauczycielami. \\
Moderator kategorii & Nauczyciel przypisany przez administratora do kategorii; akceptuje kursy, komentarze i~ocenia testy otwarte w~swojej gałęzi drzewa. \\
Reputacja / ranga & Suma głosów ,,lubię / nie lubię'' oddanych na komentarze użytkownika i~wynikający z~niej tytuł (Obserwator, Nowicjusz, Pomocnik, Ekspert, Mistrz). \\
Motyw & Nazwany zestaw pięciu kolorów bazowych (\code{ColorTheme}); odcienie pochodne i~kolory tekstu są wyliczane automatycznie. \\
Układ slajdera & Nazwana konfiguracja kart w~galerii (\code{SliderPreset}); tylko jeden jest wyświetlany publicznie. \\
Kod dostępu & Jednorazowy kod wydawany za drugą i~kolejną sztukę tego samego kursu w~zamówieniu; pozwala zapisać na kurs inną osobę. \\
Styl uczenia się & Preferowany sposób przyswajania wiedzy (\code{LearningStyle}); domyślnie cztery style VARK: wzrokowiec, słuchowiec, czytanie i~pisanie, praktyk. Administrator dodaje kolejne. \\
Ankieta stylu & Dziesięć pytań sytuacyjnych; każda odpowiedź wskazuje jeden styl. Wynikiem jest procentowy rozkład stylów i styl dominujący. \\
Profil stylów ucznia & Udział każdego stylu u ucznia wyliczony z wybranego stylu, wyniku ankiety i stylu, w którym uczniowi najlepiej idą testy; podstawa polecania kursów. \\
Wersja lekcji & Treść lekcji przygotowana dla jednego stylu (\code{LessonVariant}); uczeń bez własnej wersji widzi treść podstawową. \\
Asystent, Bielik & Czat z~polskim modelem językowym Bielik (Ollama), znający kontekst bieżącej lekcji lub wyniku testu. \\
JWT & JSON Web Token -- token dostępu używany przez API czatu i~koszyka. \\
\bottomrule
\end{tabularx}
\end{table}

\section{Zakres systemu}

Zakres funkcjonalny platformy podzielono na dziewięć modułów, odpowiadających zarazem
aplikacjom Django w~katalogu \code{apps/}.

\begin{enumerate}[label=\textbf{M\arabic*.},leftmargin=1.2cm]
  \item \textbf{Konta i~role} -- rejestracja, logowanie hasłem oraz przez Discord, GitLab
        i~Bitbucket (OAuth2), zmiana hasła, czasowa blokada po pięciu nieudanych logowaniach,
        profil z~adresem i~awatarem, trzy role, tokeny JWT, statystyki platformy dla administratora.
  \item \textbf{Katalog} -- drzewo kategorii z~okruszkami nawigacyjnymi, wyszukiwarka,
        filtry (kategoria, poziom, cena), sortowanie (także ,,dopasowane do mnie''), sekcja
        ,,Polecane dla Ciebie'', szczegóły kursu z~liczbą osób, które kurs zaczęły i~ukończyły.
  \item \textbf{Tworzenie i~moderacja treści} -- trzyetapowy kreator kursu (dane, program,
        publikacja), zgłoszenia z~wątkiem dyskusji, kolejki moderacji, zarządzanie
        kategoriami i~przydziałem moderatorów.
  \item \textbf{Społeczność} -- komentarze i~recenzje w~wątkach z~akceptacją, edycja,
        głosowanie, reputacja i~rangi, centrum powiadomień.
  \item \textbf{Sklep} -- koszyk (także dla gościa) z~kumulacją ilości, zamówienie z~adresem
        z~profilu, e-mail potwierdzający, płatność Stripe w~trybie testowym z~webhookiem,
        kody dostępu, obsługa statusów zamówień.
  \item \textbf{Nauka} -- lekcje z~treścią sformatowaną, filmem i~załącznikami, postęp,
        testy ABC i~otwarte (tekst, kod, zdjęcie, plik), ocenianie ręczne, dyplomy PDF
        z~publiczną weryfikacją, style uczenia się: wybór stylu przez ucznia, wersje lekcji
        dla stylów, ankieta diagnostyczna, rekomendacja stylu na podstawie wyników testów
        oraz polecanie kursów dopasowanych do profilu stylów ucznia; panel postępów uczniów
        kursu dla nauczyciela.
  \item \textbf{Asystent AI} -- okno czatu na każdej stronie, strumieniowanie odpowiedzi,
        kontekst lekcji i~wyniku testu, blokada w~trakcie rozwiązywania testu, historia
        rozmów.
  \item \textbf{Galeria} -- dodawanie zdjęć ze skalowaniem, konfigurator układów slajdera
        z~przeciąganiem, zapętlony slajder publiczny.
  \item \textbf{Wygląd i~dostępność} -- motywy kolorystyczne, autorski system stylów,
        pełna obsługa klawiaturą i~zgodność z~WCAG 2.1 AA.
\end{enumerate}

\subsection*{Poza zakresem}

Świadomie \emph{nie} realizowano: płatności produkcyjnych (Stripe działa wyłącznie
w~trybie testowym), wystawiania faktur, transmisji wideo na żywo, aplikacji mobilnej
(serwis jest w~pełni responsywny), hostingu modelu językowego w~chmurze oraz
wielojęzyczności interfejsu (UI jest po polsku).

\section{Założenia projektowe}

\begin{itemize}
  \item \textbf{Jeden monolit, wiele aplikacji.} Każdy moduł to osobna aplikacja Django
        z~własnymi modelami, widokami, adresami i~testami; współdzielą one użytkownika,
        kurs i~system powiadomień.
  \item \textbf{Renderowanie po stronie serwera.} Strony generują szablony Django;
        JavaScript jedynie \emph{ulepsza} interfejs (strumień czatu, zmiana ilości w~koszyku
        bez przeładowania, przeciąganie kart). Każda akcja działa też bez skryptów, przez
        zwykły formularz.
  \item \textbf{Zmiany stanu wyłącznie metodą POST} z~tokenem CSRF; jedynym wyjątkiem jest
        webhook płatności weryfikowany podpisem HMAC.
  \item \textbf{Uprawnienia sprawdzane w~widoku}, a~nie tylko przez ukrywanie przycisków.
  \item \textbf{Bez zewnętrznych usług dla danych użytkownika}: zdjęcia są skalowane
        lokalnie (Pillow), model AI działa na własnym serwerze, a~wypowiedzi
        uczniów nie trafiają do chmury.
\end{itemize}
